\chapter*{GLOSSÁRIO}
\addcontentsline{toc}{chapter}{GLOSSÁRIO}
\label{chap:glossario}

\begin{description}

\item[\textbf{Action (ROS~2)}]
Protocolo de comunicação assíncrono que permite enviar uma meta (\textit{goal})
a um servidor, receber atualizações periódicas de progresso (\textit{feedback})
e um resultado final (\textit{result}). Adequado para operações de longa duração,
como navegação autônoma. Difere de serviços (síncronos) e tópicos (sem garantia
de resultado).

\item[\textbf{AMCL} (\textit{Adaptive Monte Carlo Localization})]
Algoritmo de localização probabilística baseado em filtro de partículas que estima
a pose do robô em um mapa pré-construído. Cada partícula representa uma hipótese
de pose; as partículas são reponderadas e reamostradas a cada observação do LiDAR.
Diferente do SLAM, o AMCL pressupõe que o mapa é fixo e conhecido.

\item[\textbf{Artefato de Pesquisa} (\textit{Research Artifact})]
Código, dados, configurações e demais materiais que, publicados junto a um
estudo empírico, permitem que terceiros reproduzam ou repliquem seus
resultados. A política de \textit{badges} da ACM (\textit{Artifacts
Available}, \textit{Evaluated -- Functional}, \textit{Results Replicated})
formaliza níveis crescentes de rigor nesse processo \cite{guevaravega2024}.

\item[\textbf{Bag (rosbag2 / MCAP)}]
Arquivo de gravação de mensagens ROS~2 ao longo do tempo. O formato MCAP
(\textit{Multi-Channel Autonomous Perception}) é o padrão a partir do ROS~2 Humble,
oferecendo indexação eficiente, busca por timestamp e suporte a múltiplos tipos de
mensagem. Bags são a fonte primária de dados para análise offline de experimentos.

\item[\textbf{Behavior Tree (BT)}]
Estrutura de controle hierárquica que organiza comportamentos de um agente autônomo
em nós de ação, condição, sequência e seleção. O Nav2 utiliza BTs para orquestrar
o planejamento de trajetória, controle local e recuperação de falhas. Definidos em
arquivos XML, os BTs são facilmente modificáveis sem recompilação.

\item[\textbf{BEST\_EFFORT (QoS)}]
Política de QoS que prioriza latência mínima em detrimento de garantias de entrega.
Adequada para sensores de alta frequência (LiDAR, IMU) onde uma mensagem perdida
é aceitável se a próxima chegar em seguida. Incompatível com subscribers RELIABLE.

\item[\textbf{CDR} (\textit{Common Data Representation})]
Formato binário de serialização usado pelo DDS (e, por extensão, pelo ROS~2)
para codificar mensagens trocadas entre nós e armazenadas em \textit{bags}.
A leitura de um \textit{bag} sem o pacote de mensagens correspondente
instalado exige desserialização manual do CDR contra o tipo esperado.

\item[\textbf{Dead Reckoning}]
Estimativa de pose de um veículo por integração de velocidade e direção a partir
de uma pose inicial conhecida, sem observações externas. Em robótica móvel, é
realizado via odometria de roda. Acumula erro ao longo do tempo.

\item[\textbf{DDS} (\textit{Data Distribution Service})]
Middleware de comunicação publicar-subscrever que serve de base ao ROS~2.
Define políticas de QoS, descoberta automática de participantes na rede e
comunicação ponto-a-ponto sem broker central.

\item[\textbf{DiffDrive} (\textit{Differential Drive})]
Plugin do Gazebo que simula o acionamento diferencial de um robô com duas rodas
independentes. Publica odometria e transforma comandos de velocidade linear/angular
em velocidades individuais de roda via cinemática inversa.

\item[\textbf{Endpoint Error}]
Ver \textbf{Erro de Endpoint}.

\item[\textbf{Erro de Endpoint}]
Distância euclidiana entre as posições finais de duas execuções do mesmo percurso.
Complementa o RMSE ao focar no destino final em vez de toda a trajetória.
Um erro de endpoint pequeno com RMSE grande indica caminhos diferentes chegando
ao mesmo destino.

\item[\textbf{Filtro de Partículas}]
Técnica de estimação bayesiana que aproxima uma distribuição de probabilidade
--- por exemplo, a pose de um robô --- por um conjunto finito de amostras
ponderadas (partículas), reponderadas e reamostradas a cada nova observação.
Base algorítmica do AMCL.

\item[\textbf{FUUT} (\textit{Fixed Unit Under Test/Tasking})]
Papel experimental definido no framework Fleet UI. Uma unidade robótica em posição
fixa durante o experimento, responsável por observar o ambiente e coletar dados
de perspectiva externa. Não recebe comandos de navegação.

\item[\textbf{Gazebo Harmonic}]
Simulador robótico 3D com física rígida, versão \textit{harmonic} (GZ 8.x),
que é a versão de suporte estendido pareada com ROS~2 Jazzy. Successor do
Gazebo Classic, usa arquitetura modular e mensagens \texttt{gz.msgs}.

\item[\textbf{Goal Handle}]
Objeto retornado por um cliente de \textit{action} ROS~2 ao ter uma meta
aceita por um servidor, usado para consultar o status da execução, obter o
resultado final de forma assíncrona ou solicitar cancelamento. Ver também
\textbf{Action (ROS~2)}.

\item[\textbf{IMU} (\textit{Inertial Measurement Unit})]
Sensor que mede aceleração linear (acelerômetro) e velocidade angular (giroscópio)
em três eixos. Utilizado para estimativa de atitude e fusão com odometria de roda
para redução de deriva.

\item[\textbf{Lifecycle (nó ROS~2)}]
Modelo de gerenciamento de estado para nós ROS~2 que requerem inicialização
estruturada. Estados: \textit{unconfigured} $\to$ \textit{inactive} $\to$
\textit{active} $\to$ \textit{finalized}. O Nav2 usa lifecycle extensivamente
para garantir que todos os recursos estejam alocados antes do início da navegação.

\item[\textbf{LiDAR} (\textit{Light Detection and Ranging})]
Sensor que emite pulsos de laser e mede o tempo de retorno para determinar
distâncias. O RPLIDAR S2 utilizado nesta dissertação é um sensor 2D de varredura
rotativa com alcance de até 30\,m e frequência de 10\,Hz.

\item[\textbf{MCAP}]
Ver \textbf{Bag}.

\item[\textbf{MPPI} (\textit{Model Predictive Path Integral Control})]
Algoritmo de controle ótimo por amostragem. A cada ciclo, gera $K$ trajetórias
candidatas com ruído gaussiano, avalia seus custos, e retorna a média ponderada
como controle ótimo. É o controlador local padrão do Nav2 na versão Jazzy.

\item[\textbf{MUUT} (\textit{Mobile Unit Under Test/Tasking})]
Papel experimental definido no framework Fleet UI. A unidade robótica que executa
e grava percursos experimentais. Recebe comandos de navegação do orquestrador.
Pode haver múltiplas unidades MUUT por experimento.

\item[\textbf{Nav2}]
Navigation2 --- a pilha de navegação autônoma padrão para ROS~2, introduzida
pelo \textit{Marathon 2} \cite{macenski2020marathon}. Composta por servidores
de planejamento global, controle local, comportamentos de recuperação e um
orquestrador por Behavior Trees.

\item[\textbf{OccupancyGrid}]
Representação de mapa 2D como grade de células, onde cada célula armazena a
probabilidade de ocupação (0--100, onde -1 indica desconhecido). Formato padrão
em ROS~2 para mapas de navegação.

\item[\textbf{Odometria}]
Estimativa de pose de um robô derivada de sensores de movimento internos
(encoders de roda, IMU). Publicada como \texttt{nav\_\allowbreak msgs/\allowbreak msg/\allowbreak Odometry} no ROS~2.

\item[\textbf{QoS} (\textit{Quality of Service})]
Conjunto de políticas que controlam o comportamento da comunicação entre nós
ROS~2, incluindo confiabilidade (\textit{reliability}), durabilidade e histórico
de mensagens. Incompatibilidades de QoS entre publicador e subscritor são
silenciosas --- nenhuma mensagem é entregue sem aviso de erro.

\item[\textbf{Record--Replay}]
Protocolo experimental do framework Fleet UI composto por duas fases:
\textit{record} (gravação do percurso baseline por navegação sequencial de waypoints)
e \textit{replay} (reprodução automática do percurso via \texttt{FollowWaypoints}).

\item[\textbf{RMSE} (\textit{Root Mean Square Error})]
Métrica de distância média entre dois conjuntos de pontos, calculada como a raiz
da média dos quadrados das distâncias euclidianas ponto a ponto. Amplamente
utilizada para avaliar repetibilidade de trajetórias em robótica experimental.

\item[\textbf{ROS~2}]
Robot Operating System 2 --- middleware de código aberto para desenvolvimento
de aplicações robóticas, construído sobre o protocolo DDS. Versão Jazzy Jalisco
(LTS) utilizada nesta dissertação.

\item[\textbf{SLAM} (\textit{Simultaneous Localization and Mapping})]
Processo pelo qual um robô constrói um mapa do ambiente e, simultaneamente,
estima sua posição dentro desse mapa, sem conhecimento prévio de nenhum dos dois.
Ver também \textbf{SLAM Toolbox}.

\item[\textbf{SLAM Toolbox}]
Implementação de SLAM 2D baseada em grafos de poses para ROS~2
\cite{macenski2021slam}. Suporta mapeamento lifelong e modos síncrono/assíncrono.
É a implementação recomendada para uso com Nav2.

\item[\textbf{SU} (\textit{Support Unit})]
Papel experimental definido no framework Fleet UI. Unidade de infraestrutura
que monitora a rede ou apoia a execução do experimento sem participar diretamente
da navegação ou coleta de dados.

\item[\textbf{TF2} (\textit{Transform Library 2})]
Biblioteca do ROS~2 para gerenciamento de transformações geométricas entre
referenciais. Mantém um grafo temporal de transformações que permite calcular a
posição de qualquer referencial em relação a qualquer outro em qualquer instante
de tempo dentro do buffer de cache.

\item[\textbf{TRANSIENT\_LOCAL (QoS)}]
Política de durabilidade que mantém a última mensagem publicada em cache, entregando-a
imediatamente a novos subscritores. Necessária para tópicos como \texttt{/\allowbreak amcl\_\allowbreak pose},
onde o estado mais recente deve estar disponível para nós que ingressam após a publicação.

\item[\textbf{Waypoint}]
Ponto intermediário ou destino em uma rota de navegação, especificado como pose
$(x, y, \theta)$ no referencial do mapa. O Nav2 aceita listas de waypoints via
a ação \texttt{FollowWaypoints}.

\end{description}
